% answers_short.tex —— 2018 年·全国Ⅲ卷·理科 · 简短答案速查
% 由 tools/make_short_answers.py 自 answers.tex 的 \gksolhead/\gkanswer 提取，
% 并与 2018/2018_数学真题_全国3理_答案要点.md 逐题比对后生成，勿手改（改源 MD 后重跑）。

\begin{gknotice}
\textbf{使用说明}：本节只给最终答案便于快速核对，完整推导见解析版。答案为编写者独立推导，\textbf{非官方标准答案}。
\end{gknotice}

\gkpart{选择题}
\begin{center}\begin{tabular}{ccccccccccccc}
\toprule
题号 & 1 & 2 & 3 & 4 & 5 & 6 & 7 & 8 & 9 & 10 & 11 & 12 \\
\midrule
答案 & C & D & A & B & C & A & D & B & C & B & C & B \\
\bottomrule
\end{tabular}\end{center}

\gkpart{填空题}
\begin{itemize}[leftmargin=2.6em]
  \item[\textbf{13.}] $\dfrac12$
  \item[\textbf{14.}] $-3$
  \item[\textbf{15.}] $3$
  \item[\textbf{16.}] $2$
\end{itemize}

\gkpart{解答题}
\begin{itemize}[leftmargin=2.6em]
  \item[\textbf{17.}] (1) $a_n = 2^{n-1}$ 或 $a_n = (-2)^{n-1}$；(2) $m = 6$
  \item[\textbf{18.}] (1) 第二种生产方式效率更高；(2) 中位数 $m=80$；按超过、不超过 $m$ 的顺序，第一种分别为 15、5 人，第二种分别为 5、15 人；(3) 计算得 $K^2 = 10$。因为 $K^2 = 10 > 6.635$，所以有 99\% 的把握认为两种生产方式的效率有差异
  \item[\textbf{19.}] (1) 证明见解析版；(2) $\dfrac{2\sqrt5}{5}$
  \item[\textbf{20.}] (1) 证明见解析版；(2) 证明见解析版；公差为 $-\dfrac{3\sqrt{21}}{28}$ 或 $\dfrac{3\sqrt{21}}{28}$
  \item[\textbf{21.}] (1) 证明见解析版；(2) $a = -\dfrac16$
\end{itemize}

\gkpart{选做题}
\begin{itemize}[leftmargin=2.6em]
  \item[\textbf{22.}] (1) $\alpha \in \left(\dfrac\pi4, \dfrac{3\pi}{4}\right)$；(2) $\begin{cases} x = \sqrt2\sin\alpha\cos\alpha, \\ y = -\sqrt2 + \sqrt2\sin^2\alpha \end{cases} \quad \left(\alpha \in \left(\dfrac\pi4, \dfrac{3\pi}{4}\right)\right)$
  \item[\textbf{23.}] (1) $f(x) = \begin{cases} -3x, & x < -\dfrac12, \\ x + 2, & -\dfrac12 \leqslant x \leqslant 1, \\ 3x, & x > 1\text{。} \end{cases}$ 图象为连接点 $\left(-\dfrac12, \dfrac32\right)$、$(1, 3)$ 的分段折线图象；(2) $a + b$ 的最小值为 $5$
\end{itemize}
